\documentclass[10pt,a4paper]{amsart}
\usepackage[margin=19mm]{geometry}
\usepackage{amsmath,amssymb,mathtools}
\usepackage[scr=rsfs,cal=euler]{mathalfa}
\usepackage{xfrac}
\usepackage[unicode,hidelinks,bookmarks=false]{hyperref}
\newcommand{\ie}{i.e.\ }
\newcommand{\cf}{cf.\ }
\newcommand{\resp}{respectively\ }
\newcommand{\Subsection}{\subsection}
\providecommand{\nobreakdash}{\nobreak\hbox{-}}
\setlength{\emergencystretch}{2em}
\multlinegap0pt
\usepackage{amssymb}
\usepackage{bm}
\newtheorem{thm}{Theorem}
\title{\textbf{Spectral bounds for the independence number of graphs and even uniform hypergraphs}}
\author{Xinyu Hu, Jiang Zhou, Changjiang Bu}
\date{Source: arXiv:2603.07501v1 (CC BY 4.0)}
\begin{document}
\maketitle

\noindent\emph{Source and licence.} \textbf{Spectral bounds for the independence number of graphs and even uniform hypergraphs}. Version arXiv:2603.07501v1 (CC BY 4.0), distributed by arXiv under the Creative Commons Attribution 4.0 International licence, http://creativecommons.org/licenses/by/4.0/, as recorded in the arXiv licence metadata for this exact version and shown on the abstract page https://arxiv.org/abs/2603.07501v1. Full licence notice: \texttt{licenses/arxiv-2603.07501-CC-BY-4.0.txt}.\par\smallskip

\noindent\textit{Editorial scope.} Counted unit: the article's thm t (source0.tex lines 201--210). The article's complete original proof of it is delivered with it (source0.tex lines 211--255, one \texttt{proof} environment of 45 lines). No mathematics is written, repaired or re-derived: the statement and the proof are the article's own lines, in the article's own notation, and no retained line is substituted or re-flowed.  No further source-local context is needed: every cross-reference of every retained line resolves inside the two delivered spans.  Nothing was omitted from any retained span, so the package carries no omission record.  No retained line contains a citation, so the wrapper carries no bibliography and no bibliography entry.  No retained line contains a figure environment, an image-inclusion command or drawing code, so no image is delivered and the package declares no asset dependency.  Editorial formatting only, in addition: an AMS \texttt{amsart} preamble that restates the article's own declarations and macros used by the retained lines, namely the environments \texttt{thm} on separate counters, together with the standard packages those lines need.  The retained environments print in this document as Theorem 1 in order of appearance, whereas the article prints the same items with the numbering of its own full document.  The complete native source of the article as distributed in the arXiv e-print for this exact version is bundled unchanged as \texttt{sources/195949/source0.tex}.
\begin{thm}\label{t}
Suppose that $t$ is an odd integer such that $0<t<k$ for an even integer $k$. Let $\mathcal{H}$ be a $k$-uniform hypergraph with $n$ vertices and $m\geq1$ edges, and let $\delta$ and $\lambda$ be the minimum degree and the minimum $H$-eigenvalue of $\mathcal{H}$, respectively. Then
\begin{eqnarray*}
{\alpha _t}(\mathcal{H}) \leqslant \frac{{t\left( {km - n \lambda } \right){{\left( { - \lambda } \right)}^{\frac{t}{{k - t}}}}}}{{\left( {k - t} \right){\delta ^{\frac{k}{{k - t}}}} + \left( {k \delta  - t\lambda } \right){{\left( { - \lambda } \right)}^{\frac{t}{{k - t}}}}}},
\end{eqnarray*}
with equality if and only if there exits a $t$-independent set $S$ such that every vertex in $S$ has degree $\delta$, and every vertex $i$ outside $S$ satisfies
\begin{eqnarray*}
\left| {\{ e:e \in E(i),e\cap S\neq\emptyset \} } \right| = \frac{{{{\left( { - \lambda } \right)}^{\frac{k}{{k - t}}}} + {d_i}{{\left( { - \lambda } \right)}^{\frac{t}{{k - t}}}}}}{{{\delta ^{\frac{t}{{k - t}}}} + {{\left( { - \lambda } \right)}^{\frac{t}{{k - t}}}}}}.
\end{eqnarray*}
\end{thm}

\begin{proof}
Since $m\geq1$, we have $\lambda<0$. Let $S$ be a $t$-independent set of $\mathcal{H}$ such that $|S|=\alpha_t(\mathcal{H})$. For $c\geq0$, take $x=(x_1,x_2,\ldots,x_n)^\top$ satisfying
$${x_i} =
\begin{cases}
c & i \in S,\\
- 1 & i \notin S.
\end{cases}
$$
Let $\mathcal{A}$ be the adjacency tensor of $\mathcal{H}$. Since $\mathcal{A}-\lambda \mathcal{I}$ is a positive semidefinite tensor, we have
\begin{eqnarray*}
(\mathcal{A}-\lambda \mathcal{I})x^k=-\lambda\sum_{i=1}^nx_i^{k}+k\sum_{\{ i_1,i_2,\ldots,i_k\}\in E}x_{i_1}x_{i_2}\cdots x_{i_k}\geq0,
\end{eqnarray*}
where $E=E(\mathcal{H})$ is the edge set of $\mathcal{H}$.
Since $|S|=\alpha_t(\mathcal{H})$ and $|e\cap S|\in\{0,t\}$ for each edge $e\in E$, we have

\[( {\mathcal{A}}-\lambda \mathcal{I} )x^k = -\lambda\sum_{i=1}^nx_i^k  +k
 \sum_{\substack{\{i_1,i_2,\ldots,i_k\}\in E\\ |\{i_1,i_2,\ldots,i_k\}\cap S|=t}}
{{x_{{i_1}}}{x_{{i_2}}} \cdots {x_{{i_k}}}}  +k
\sum_{\substack{\{i_1,i_2,\ldots,i_k\} \in E\\ \{i_1,i_2,\ldots,i_k\} \cap S = \emptyset}} {{x_{{i_1}}}{x_{{i_2}}} \cdots {x_{{i_k}}}} \]


\[= - {\lambda} \alpha_{t}(\mathcal{H}) {c^k} - \lambda (n - {\alpha_{t}}(\mathcal{H}) ) - kc^tt^{-1}\sum_{i\in S}d_i+ k\left( {m - t^{-1}\sum_{i\in S}d_i}\right) \geqslant 0.\]
Then
\[{\alpha _t}(\mathcal{H})\left( {-\lambda {c^k} + \lambda } \right) + km -n \lambda  \geqslant \left( {k{c^t} + k} \right)t^{-1}\sum_{i\in S}d_i \geqslant \left( {k{c^t} + k} \right)\frac{{{\alpha _t}(\mathcal{H})\delta }}{t},\]

\[\left( {-t \lambda \left( {1 - {c^k}} \right) + k \delta \left( {1 + {c^t}} \right)} \right){\alpha _t}(\mathcal{H})\leq t\left( {km - n\lambda } \right).\]
Take $c = {{\left( \frac{\delta }{-\lambda } \right)}^{\frac{1}{k-t}}}$, then
$ -t \lambda \left( {1 - {c}^k} \right) + k \delta \left( {1 + {c}^t} \right) = \left( {k - t} \right)\delta {c}^t + k \delta - t \lambda  > 0$ and
\[{\alpha _t}(\mathcal{H}) \leqslant \frac{{t\left( {km - n \lambda } \right)}}{{\left( {k - t} \right)\delta {c^t} + k \delta  - t \lambda }}=\frac{{t\left( {km - n \lambda } \right){{\left( { - \lambda } \right)}^{\frac{t}{{k - t}}}}}}{{\left( {k - t} \right){\delta ^{\frac{k}{{k - t}}}} + \left( {k \delta  - t\lambda } \right){{\left( { - \lambda } \right)}^{\frac{t}{{k - t}}}}}},\]
with equality if and only if every vertex in $S$ has degree $\delta$, and $x$ is an eigenvector of $\lambda$. Recall that $x$ is a vector satisfying
$${x_i} =
\begin{cases}
 {{\left( \frac{\delta }{-\lambda } \right)}^{\frac{1}{k-t}}} & i \in S,\\
-1 & i \notin S.
\end{cases}
$$
When every vertex in $S$ has degree $\delta$, $\mathcal{A}x^{k-1}=\lambda x^{[k-1]}$ is equivalent to
\begin{eqnarray*}
-\lambda=\lambda x_i^{k-1}=\sum_{\{i,i_2,\ldots,i_k\}\in E}x_{i_2}\cdots x_{i_k},~i\notin S.
\end{eqnarray*}
For $i\notin S$, let $d'_i = \left| \{e: e \in E(i),e\cap S \neq \emptyset \}\right|$. Then the above equation is equivalent to
$$-\lambda = d'_i {{\left( \frac{\delta }{-\lambda } \right)}^{\frac{t}{k-t}}} - (d_i - d'_i),$$
that is,
$$d'_i = \left| {\{ e:e \in E(i),e\cap S \ne \emptyset \} } \right| = \frac{{{{\left( { - \lambda } \right)}^{\frac{k}{{k - t}}}} + {d_i}{{\left( { - \lambda } \right)}^{\frac{t}{{k - t}}}}}}{{{\delta ^{\frac{t}{{k - t}}}} + {{\left( { - \lambda } \right)}^{\frac{t}{{k - t}}}}}}.$$
\end{proof}

\end{document}
